% Generated draft: requires mathematical and source review.
\documentclass[11pt]{article}
\usepackage[margin=25mm]{geometry}
\usepackage{fontspec}
\setmainfont{Latin Modern Roman}
\usepackage{amsmath,amssymb,mathtools,unicode-math}
\setmathfont{Latin Modern Math}
\usepackage{enumitem,needspace,xurl,hyperref}
\setlength{\parindent}{0pt}
\setlength{\parskip}{5pt}
\setlength{\emergencystretch}{3em}
\newcommand{\N}{\mathbb{N}}
\newcommand{\Z}{\mathbb{Z}}
\newcommand{\Q}{\mathbb{Q}}
\newcommand{\R}{\mathbb{R}}
\newcommand{\C}{\mathbb{C}}
\newcommand{\F}{\mathbb{F}}
\newcommand{\A}{\mathbb{A}}
\newcommand{\PP}{\mathbb{P}}
\newcommand{\E}{\mathbb{E}}
\newcommand{\eps}{\varepsilon}
\newcommand{\dd}{\,\mathrm{d}}
\begin{document}
\section*{UnsolvedMath ID: 30006421}
\textbf{Source problem code:} catalogue entry
\subsection{Definitions and mathematical statement}
Let \((M^4,g)\) be a smooth, complete, contractible Riemannian four-manifold. Contractibility means that the identity map of \(M\) is homotopic to a constant map. Write \(d_g\) for the Riemannian distance and \(\operatorname{Vol}_g\) for Riemannian volume. Nonnegative Ricci curvature means
\[
\operatorname{Ric}_g(v,v)\geq 0\qquad\text{for every }x\in M\text{ and }v\in T_xM.
\]
Use the conventional meaning of Euclidean volume growth:
\[
\liminf_{r\to\infty}\frac{\operatorname{Vol}_g(B_g(p,r))}{r^4}>0
\]
for some base point \(p\in M\). Here \(B_g(p,r)\) is the metric ball of radius \(r\). This growth condition is independent of the choice of the fixed base point.

\textbf{Conjecture 1 (as supplied).} Under these hypotheses, both of the following assertions hold:
\begin{enumerate}
\item \(M\) is homeomorphic to \(\R^4\).
\item For every \((Z,d_Z)\in\mathcal C_\infty\), the space \(Z\) is homeomorphic to
\[
S^3=\{x\in\R^4:\|x\|=1\}.
\]
\end{enumerate}
These are assertions of homeomorphism; no diffeomorphism conclusion is stated.

\textbf{Unresolved definition of the second assertion.} The supplied original passage uses \(\mathcal C_\infty\) without defining it, and the cleaned statement calls its members asymptotic cross-sections. A conventional interpretation is the collection of cross-sections of all metric cones arising as tangent cones at infinity of \((M,d_g)\): a compact metric space \((Z,d_Z)\) belongs to this collection when, for some sequence \(r_i\to\infty\),
\[
(M,r_i^{-1}d_g,p)\longrightarrow (C(Z),d_C,o)
\]
in the pointed Gromov--Hausdorff sense. Here \(C(Z)\) is the metric cone over \(Z\), and \(o\) is its vertex. Under this interpretation, assertion (ii) quantifies over every such cross-section and every such limiting sequence, without assuming uniqueness of the tangent cone. The identification of this conventional collection with the source's \(\mathcal C_\infty\) requires confirmation from the surrounding report; it is an explicit interpretation, not a definition supplied by the excerpt.
\subsection{Short English statement}
The supplied conjecture asks whether every smooth, complete, contractible four-manifold with nonnegative Ricci curvature and Euclidean volume growth is homeomorphic to \(\R^4\), and whether every asymptotic cross-section is homeomorphic to \(S^3\). The source notation for the collection of cross-sections is undefined in the supplied excerpt, so the second assertion requires confirmation of its intended scope.
\subsection{Sources}
\begin{itemize}
\item Supplied local UnsolvedMath registry record, permanent ID \(30006421\), source problem code \texttt{catalogue entry}, titled ``Contractible Ricci-Nonnegative Four-Manifolds with Euclidean Growth,'' in the Oberwolfach Reports Open Problems collection. The original excerpt labels the statement ``Conjecture 1.'' No individual proposer is supplied; the proposed year is recorded as 2025.
\item Supplied report citation: \emph{Differentialgeometrie im Großen (2026)}, Oberwolfach Reports, pp.\ 1591--1646, DOI \texttt{10.4171/owr/2025/30}. \url{https://doi.org/10.4171/owr/2025/30}. The citation's 2026 title date and the DOI's 2025 component are retained as supplied.
\item Supplied literature reference: Brué--Pigati--Semola survey. \url{https://doi.org/10.1016/j.indag.2025.04.003}.
\item Supplied literature reference: Huang--Huang fundamental groups. \url{https://arxiv.org/abs/2502.03259}.
\item Supplied literature reference: Oberwolfach Report 30. \url{https://ems.press/content/serial-article-files/52249}.
\item Literature assessment dated \texttt{2026-08-22}: the supplied curation labels the conjecture open, says that a 2026 survey still asks about the existence of a nonstandard contractible four-manifold with nonnegative Ricci curvature and Euclidean volume growth, and describes recent work as providing structural and fundamental-group restrictions only. This is supplied dated curation derived from the OWR record, not an independently verified literature assessment. The record explicitly requests independent verification of time-sensitive status and cited claims.
\end{itemize}
\end{document}
